\section{An obstruction to a uniform critical-rate counting theorem}
\label{sec:critical-rate-obstruction}

There is a stronger reason to avoid treating the critical exponent as
only a technical extension of the preceding proof. Absolute critical-rate
sublevel sets cannot have a counting bound jointly polynomial in input
dimension, depth, and rate. This is a statement about the actual
mean-field rate functional, although it is still not a sample-complexity
lower bound for a simple target.

We use the critical finite-domain variational formula proved in
Appendix~\ref{sec:critical-variational}.
For a feature law $\nu$ on $[-1,1]^{\mathcal X}$, let
$Q(\nu)=\mathbb E_\nu hh^T$, and let $\mathsf M_Q$ be the law of
coordinatewise $\tanh(G)$ for $G\sim\mathcal N(0,Q)$.
The critical rate of a global-margin Boolean function is
\[
 J_L(f)=\inf\left\{
 D(\nu_1\Vert\mathsf M_{Q_0})+
 \sum_{s=2}^{L-1}D(\nu_s\Vert\mathsf M_{Q(\nu_{s-1})})+
 \frac12z^TQ(\nu_{L-1})^+z:
 f(x)z(x)\ge\frac12\ \forall x\right\},
\]
where the final term is infinite outside the covariance range and
$Q_0(x,y)=\bar x\cdot\bar y/d$.

\begin{theorem}[Every Boolean function has polynomial critical rate at a polynomial depth]
\label{thm:critical-counting-obstruction}
For each $d\ge2$, set
\[
 r=20(d+1)^5,\qquad L=r+4.
\]
Every Boolean function on $\{-1,1\}^d$ satisfies
\[
 J_L(f)\le100000L.
\]
Consequently no polynomial $P$ can satisfy
\[
 \log\#\{f:J_L(f)\le K\}\le P(d,L,K)
\]
for every $d,L,K$.
\end{theorem}

\begin{proof}
Write $m=2^d$. First we obtain a crude quantitative lower bound on the
smallest eigenvalue of the unperturbed first hidden covariance
\[
 Q=\mathbb E_w h(w)h(w)^T,\qquad
 h(w)_x=\tanh(w\cdot\bar x),\qquad
 w\sim\mathcal N(0,I_{d+1}/d).
\]
For each $a\in\{-1,1\}^d$, choose the center
\[
 w_a=\bigl(-b(d-1/2),ba_1,\ldots,ba_d\bigr),
 \qquad b=\log(16m).
\]
At $x=a$ its preactivation is $b/2$, whereas at every other cube point
it is at most $-3b/2$. Thus its feature vector is within $1/(8m)$
coordinatewise of the vector that is $+1$ at $a$ and $-1$ elsewhere.
Let $E_a$ be the coordinate box of radius
$\rho=1/[8m(d+1)]$ around $w_a$. These boxes are disjoint, and changing
weights within a box changes every activation by at most $1/(8m)$.
The conditional mean vectors
$\bar h_a=\mathbb E[h(w)\mid w\in E_a]$ therefore form a matrix
$H$ within operator norm $1/4$ of $2I-\boldsymbol1\boldsymbol1^T$.
Since $m\ge4$, that matrix has smallest singular value at least two;
hence $HH^T\succeq I$.

Every box has Gaussian probability at least
\[
 (2\rho)^{d+1}\left(\frac d{2\pi}\right)^{(d+1)/2}
 \exp\left[-\frac d2
       (\|w_a\|+\sqrt{d+1}\rho)^2\right]
 \ge p:=e^{-20(d+1)^5}.
\]
For the last crude estimate, use
$b\le3(d+1)$, $\|w_a\|+\sqrt{d+1}\rho\le4(d+1)^2$,
$(d+1)\log(1/(2\rho))\le3(d+1)^2$, and
$(d/(2\pi))^{(d+1)/2}\ge e^{-(d+1)}$.
Conditional Jensen's inequality, together with disjointness, gives
\[
 Q\succeq\sum_a\Pr(E_a)\bar h_a\bar h_a^T\succeq pI.
\]

Fix any Boolean vector $f\in\{-1,1\}^m$ and put
$\eta=\sqrt{p/m}\le1$. Take the first feature law to be its
unperturbed law, at zero entropy cost. For the second feature law,
shift its Gaussian preactivation mean to $\eta f$, retaining covariance
$Q$. Gaussian relative entropy and data processing give cost at most
\[
 \frac12\eta^2 f^TQ^{-1}f\le\frac12.
\]
All diagonal entries of $Q$ coincide, so the new feature mean is
exactly $m_0f$ for some $m_0>0$. It satisfies $m_0\ge\eta/32$:
for $0\le t\le1$ and variance $q\le1$,
\[
 \mathbb E\operatorname{sech}^2(t+\sqrt qZ)
 \ge\Pr\{|Z|\le1\}\operatorname{sech}^2(2)>1/32,
\]
and integration from zero to $\eta$ proves the claim.

For any feature law with mean vector $u$, its second moment obeys
$Q\succeq uu^T$, hence $u$ lies in its range and $u^TQ^+u\le1$.
Therefore shifting the next Gaussian preactivation mean to $g u$
costs at most $g^2/2$ in entropy. All feature variances remain at most
one. In our construction all diagonal second moments remain identical,
so mean vectors remain of the form $m_sf$.
Choose $g=100$ for $r$ successive layers. If $0<m_s\le1/100$, the
preceding derivative estimate gives
\[
 m_{s+1}\ge100m_s/32>3m_s.
\]
If $m_s\ge1/100$, monotonicity gives $m_{s+1}>1/32>1/100$.
Since
\[
 \log(1/\eta)=10(d+1)^5+(d/2)\log2<11(d+1)^5,
\]
the choice $r=20(d+1)^5$ guarantees $m_r\ge1/100$.
Each amplification layer costs at most $5000$.

Use one last mean shift with $g=400$. Its mean preactivation has
signed magnitude at least four. With Gaussian variance at most one,
the expected signed tanh is greater than $2/3$, by the Gaussian tail
estimate used in Proposition~\ref{prop:critical-stochastic-padding}.
This layer costs at most $80000$.
Take the final output $z$ equal to this last feature mean. It has
signed margin greater than $2/3$, and its Gaussian readout cost is
at most $1/2$, again by $Q\succeq zz^T$.
The number of hidden layers is $1+1+r+1=r+3$, so the depth is
$L=r+4$. The total rate is at most
\[
 \tfrac12+5000r+80000+\tfrac12\le100000L.
\]
Finally, at these polynomial values of $L$ and $K=100000L$, the rate
sublevel contains every one of the $2^{2^d}$ Boolean functions. Its
logarithmic cardinality is $2^d\log2$, which exceeds every polynomial
in $d,L,K$ for sufficiently large $d$.
\end{proof}

This does not prove that a simple target is hard for the critical
posterior. Such a lower bound would require comparison with the target's
\emph{own optimal rate}, sufficiently accurately to control which
functions win after conditioning. The theorem instead rules out a
specific proposed upper-bound route: a universal counting theorem based
only on absolute critical rate, with polynomial dependence on depth.
It is compatible with more precise comparison to the target, or with
additional structural information about posterior minimizers.


\begin{proposition}[The obstruction also holds directly at finite width]
\label{prop:critical-finitewidth-capacity}
With $d,L$ as in Theorem~\ref{thm:critical-counting-obstruction}, for
all widths $N\ge\lceil e^{100(d+1)^5}\rceil$ and every Boolean $f$,
\[
 \mu_{N,1}(A_{N,f})\ge\tfrac14
                  \exp\{-200000NL\}.
\]
This is a direct statement about the original finite-width Gaussian
prior and requires no appeal to a large-deviation limit.
\end{proposition}

\begin{proof}
Retain $m=2^d$, $p=e^{-20(d+1)^5}$ and $r=20(d+1)^5$, and now put
$\eta=\sqrt{p/(4m)}$. Sample the first layer from its original law.
Let $H$ be its $N\times m$ matrix of activations and
$\widehat Q=H^TH/N$. The preceding covariance bound and entrywise
Hoeffding imply
\[
 \Pr\{\lambda_{\min}(\widehat Q)<p/2\}
 \le2m^2e^{-Np^2/(8m^2)}\le1/8.
\]
Indeed entrywise error at most $p/(2m)$ implies operator error at most
$p/2$. On the good event, set
\[
 u=\frac\eta N H\widehat Q^{-1}f;
\]
on its complement set $u=0$. Give every row of the next matrix this
mean vector $u$, preserving reference covariance $I/N$.
On the good event,
\[
 H^Tu=\eta f,\qquad
 \|u\|^2=\frac{\eta^2}{N}f^T\widehat Q^{-1}f
             \le\frac1{2N}.
\]
This changed layer has conditional relative entropy at most $N/4$.
Every conditional variance is at most one, so its signed expected
activation is at least $\eta/32$ on the good event.

At each of the next $r$ layers give every incoming weight mean $100/N$,
and at one last hidden layer use mean $400/N$, always retaining
variance $1/N$. Simultaneously over all cube inputs and all these
$r+2$ shifted hidden layers, conditional Hoeffding gives absolute
empirical-mean error at most $\eta/128$, with total failure probability
at most
\[
 2m(r+2)e^{-N\eta^2/32768}\le1/8.
\]
On this event, the first shifted empirical mean has signed value at
least $\eta/64$ everywhere. As long as its minimum signed value is
at most $1/100$, a gain-$100$ layer increases it by at least a factor
of two: its expected increase is at least a factor of three, and
$\eta/128$ is at most half the current minimum. Once the minimum
reaches $1/100$ it stays above that level, by monotonicity and the
bound $1/32-\eta/128>1/100$.
The choice of $r$ ensures $2^r\eta/64\ge1/100$.
The final gain-$400$ layer then has signed expected mean greater than
$2/3$ and signed empirical mean greater than $5/8$.

Give the final readout mean $2/N$ per entry and preserve variance
$1/N^2$. Its signed mean exceeds $5/4$, and its Gaussian noise is
uniformly at most $1/2$ except with probability
$2m e^{-N/8}\le1/8$. Its clean margin is consequently greater than
$3/4$. The proposed weight law assigns $A_{N,f}$ probability at least
$5/8$, hence at least $1/2$.

The relative entropy chain rule bounds this law's entropy from the
original Gaussian by
\[
 N/4+5000Nr+80000N+2N\le100000NL.
\]
Binary data processing as in
Proposition~\ref{prop:critical-stochastic-padding} proves the mass bound.
For completeness, the displayed width assumption dominates each of
\[
 \frac{8m^2}{p^2}\log(16m^2),\qquad
 \frac{32768}{\eta^2}\log(16m(r+2)),\qquad 8\log(16m),
\]
which are the three concentration thresholds used above. This follows
directly by substituting $m=2^d$, $p=e^{-20(d+1)^5}$ and $d\ge2$.
All adaptive means are measurable functions of preceding weights;
the inverse is used only on the event with smallest eigenvalue at
least $p/2$.
\end{proof}
